\subsection[补充]{补充\footnote{在本号中，我们使用了《交换代数》一书中尚在编写中的各章的结果，并用“交换代数”这一名称来标示它们。}}

引理 4. 设 K 是交换域，V 是 K 上的有限维向量空间，G 是由 V 的自同构组成的有限群，其阶 q 在 K 中可逆，S 是 V 的对称代数，R 是 S 中由在 G 下不变的元素组成的子代数。S 的高度为 1 的素理想 B 在 \(p = B \cap R\) 上分歧（Commutative Algebra）当且仅当存在 V 的非零元素 a 与 \(V^{*}\) 的非零元素 f，使得 B = Sa 且伪反射 \(s_{a,f}\) 属于 G。此时，分解群 \(\mathcal{G}^{\mathrm{Z}}(\mathrm{B})\) 是 G 中使 Ka 稳定的元素所成的子群，惯性群 \(\mathcal{G}^{\mathrm{T}}(\mathrm{B})\) 是 G 的循环子群 \(H_{a}\)，由 G 中以 a 为向量的诸伪反射组成。S 在 B 处的剩余域 S(B) 在 R 于 p 处的剩余域 R(p) 上是可分的，而分歧指数 \(e(\mathrm{B}/\mathfrak{p})\) ——即 B 在差积的除子 \(\text{div}(D_{S/R})\) 中的系数再加 1——等于 \(\text{Card}(H_{a})\)。

说 B 在 R 上分歧，就是指其惯性群 \(\mathcal{G}^{\mathrm{T}}(\mathrm{B})\) 不化为单位元群，换言之，存在 G 中的 \(g \neq 1\)，使得 \(g(z) \equiv z \pmod{B}\) 对所有 \(z \in S\) 都成立。由于 S 是唯一分解环，B 是主理想 Sa，且 a 必须整除所有元素 \(g(z) - z\)（\(z \in S\)）；而当 \(z \in V\) 时，这些元素都是 1 次齐次的，并且都非零（因为 \(g \neq 1\)）；因此 a 必是 1 次齐次的，换言之 \(a \in V\)。于是存在 V 上的线性型 f，使得 \(g = s_{a,f}\)。反之，若 g 是异于 1 的伪反射 \(s_{a,f}\)，则 \(g(z) \equiv z \pmod{Sa}\) 对所有 \(z \in S\) 都成立，故 g 属于素理想 B = Sa 的惯性群。这就证明了引理的第一个断言，以及 \(\mathcal{G}^{\mathrm{Z}}(\mathrm{B})\) 与 \(\mathcal{G}^{\mathrm{T}}(\mathrm{B})\) 的特征刻画。由于 q 与 K 的特征 p 互素（该特征也是 S(B) 的特征），R(p) 的扩张 S(B) 是可分的（Commutative Algebra, Chap. V, §2, no. 2, Cor. of Prop. 5）。由此即得等式 \(e(\mathrm{B}/\mathfrak{p}) = \mathrm{Card}(\mathrm{H}_{a})\)（Commutative Algebra）。由于 \(e(\mathrm{B}/\mathfrak{p})\) 与 p 互素，B 在 \(\operatorname{div}(\mathrm{D}_{\mathrm{S/R}})\) 中的系数为 \(e(\mathrm{B}/\mathfrak{p}) - 1\)（Commutative Algebra）。引理得证。

引理 5. 设 K 是交换域，S 是 K 上的分次多项式代数，R 是 S 的分次子代数。则 S 是自由分次 R-模（Algebra, Chap. II, §11, no. 2）当且仅当下列两个条件得到满足：a) R 是 K 上的分次多项式代数；

\begin{enumerate}
\def\labelenumi{\alph{enumi})}
\setcounter{enumi}{1}
\tightlist
\item
  若 \((\alpha_{1},\ldots,\alpha_{s})\) 是 K-代数 R 的一个由代数无关的齐次元素组成的生成元组，则该组是一个 S-正则序列 \(^{5}\)。
\end{enumerate}

当 S 是有限型 R-模时，b) 是 a) 的推论。

证明见 Commutative Algebra。

定理 4. 设 K 是交换域，V 是 K 上的有限维向量空间，S 是 V 的对称代数，G 是由 V 的自同构组成的有限群，R 是 S 中由在 G 下不变的元素组成的子代数。假设 q = Card G 在 K 中可逆。则下列条件等价：

\begin{enumerate}
\def\labelenumi{(\roman{enumi})}
\item
  G 由伪反射生成；
\item
  S 是自由分次 R-模；
\item
  \(\mathbf{R}\) 是 \(\mathbf{K}\) 上的分次多项式代数。
\end{enumerate}

(ii) 与 (iii) 的等价性由第 2 号及引理 5 得出。蕴涵关系 (i) \(\Longrightarrow\) (ii) 由 Th. 1 得出。

我们证明 (iii) \(\Longrightarrow\) (i)。设 \(G'\) 是 \(G\) 中由属于 \(G\) 的伪反射生成的子群，\(R'\) 是 \(S\) 中由在 \(G'\) 下不变的元素组成的子代数。我们有 \(R \subset R' \subset S\)。由引理 4，\(\text{div}(D_{S/R}) = \text{div}(D_{S/R'})\)，故 \(\text{div}(D_{R'/R}) = 0\)。于是设 \(R\) 是分次多项式代数。由于 \(R'\) 也是如此（因为 \(G'\) 由伪反射生成），引理 5 表明 \(R\)-模 \(R'\) 具有齐次基 \((Q_1, \ldots, Q_m)\)；设 \(q_i = \deg(Q_i)\)。置

\[
d = \det (\mathrm{Tr} _ {R ^ {\prime} / R} (Q _ {i} Q _ {j})), \text { cf.   Algebra,   Chap.   IX, } \S 2.
\]

\(\operatorname{div}(\mathrm{D}_{\mathrm{R}^{\prime}/\mathrm{R}})\) 为零这一事实表明 \(\operatorname{div}(d)=0\)（Commutative Algebra），即 d 属于 \(K^{*}\)。另一方面，\(\operatorname{Tr}_{\mathrm{R}^{\prime}/\mathrm{R}}(\mathrm{Q}_{i}\mathrm{Q}_{j})\) 是次数为 \(q_{i}+q_{j}\) 的齐次元，而 d 是次数为 \(2\sum_{i}q_{i}\) 的齐次元。于是 \(\sum_{i}q_{i}=0\)，即对所有 i 有 \(q_{i}=0\)，这意味着 \(R^{\prime}=R\)，从而由 Galois 理论得 \(G^{\prime}=G\)。这就证明了 G 由伪反射生成。证毕。

注记。1) 在 Th. 4 的假设下，R 的诸特征次数之积为 q（公式 (6) 与 Prop. 2 (iii)），故它们都与 K 的特征互素。这一点在第 3 号中已予断言。

\begin{enumerate}
\def\labelenumi{\arabic{enumi})}
\setcounter{enumi}{1}
\tightlist
\item
  若不假设 \(\operatorname{Card}(G)\) 在 K 中可逆，则仍有蕴涵关系 (ii) \(\Longleftrightarrow\) (iii)（cf.~引理 5）与 (ii) \(\Longrightarrow\) (i)（cf.~Exerc. 8）；但蕴涵关系 (i) \(\Longrightarrow\) (ii) 不再成立（Exerc. 9）。
\end{enumerate}

命题 6. 假设与记号同 Th. 4。设 H 是属于 G 且异于 1 的伪反射所成的集合。假设 H 生成 G。对所有 \(g \in \mathbf{G}\)，置 \(g(x) = x + f_g(x)e_g\)，其中 \(e_g \in V\)，\(f_g \in V^*\)。置

\[
\mathrm{D} = \prod_ {g \in \mathrm{H}} e _ {g} \in \mathrm{S}.
\]

\begin{enumerate}
\def\labelenumi{(\roman{enumi})}
\item
  \(S\) 在 \(R\) 上的差积是主理想 SD。
\item
  把 \(S\) 与代数 \(K[X_1, \ldots, X_l]\) 等同，其中 \((X_1, \ldots, X_l)\) 是 \(V\) 的一个基；并设 \(P_1, \ldots, P_l\) 是生成代数 \(R\) 的代数无关的齐次元素。于是 Jacobi 行列式 \(J = \det\left(\frac{\partial P_i}{\partial X_j}\right)\) 形如 \(\lambda D\)，其中 \(\lambda \in K^*\)。
\item
  \(\sum_{i=1}^{l} (\deg(P_i) - 1) = \text{Card}(H)\) 。
\item
  \(S\) 的反不变元素所成的集合是 RD。
\end{enumerate}

断言 (i) 由引理 4 得出。断言 (ii) 由 SJ 是 S 在 R 上的差积这一事实得出（Commutative Algebra）。断言 (iii) 由写出齐次多项式 D 与 J 次数相同这一事实而得到。而 (iv) 的证明则与第 4 号中给出的证明相同（Prop. 5 证明中的 b) 与 d) 两部分）。*

